\documentclass[lettersize,journal]{IEEEtran}
\usepackage{amsmath,amsfonts}
\usepackage{algorithmic}
\usepackage{algorithm}
\usepackage{array}
\usepackage[caption=false,font=normalsize,labelfont=sf,textfont=sf]{subfig}
\usepackage{textcomp}
\usepackage{stfloats}
\usepackage{url}
\usepackage{verbatim}
\usepackage{graphicx}
\usepackage{cite}
\usepackage{hyperref}
\usepackage{booktabs}
\usepackage{multirow}

\hyphenation{op-tical net-works semi-conduc-tor TopoVarAD}

\begin{document}

\title{TopoVarAD: Topology-Aware Variational Autoregressive Anomaly Detection for Weakly-Supervised High-Resolution Industrial Images}

\author{Anonymous Author,~\IEEEmembership{Author,~Affiliation}
\thanks{Manuscript received xxx; revised xxx.}}

\markboth{IEEE Transactions on xxx,~Vol.~xx, No.~xx, Month~2026}%
{Author \MakeLowercase{\textit{et al.}}: TopoVarAD: Topology-Aware Variational Autoregressive Anomaly Detection}

\maketitle

\begin{abstract}
Industrial image anomaly detection faces three fundamental challenges: diverse defect morphologies that elude rigid grid partitioning, the destruction of spatial topology by global flattening in state-space models, and the limited discriminability of continuous Gaussian-based anomaly scoring. This paper presents TopoVarAD, a topology-aware variational autoregressive framework that addresses these limitations through three coordinated innovations. First, a T2M-Tokenizer employs multi-scale SLIC superpixel segmentation with adaptive pooling and Graph Laplacian positional encoding to preserve the geometric integrity of defects within token representations. Second, a Topology-Preserving Mamba (TPM) Block combines bidirectional state-space scanning with topology-constrained sparse attention via a learnable gating mechanism, maintaining spatial structure while achieving linear complexity. Third, a Topology-Aware Autoregressive (TAR) Head replaces continuous latent space modeling with residual vector quantization and GPT-style next-token prediction, converting anomaly detection into discrete sequence likelihood estimation. The framework is trained in two stages: a reconstruction pre-training phase (Stage~1) that learns pixel-level representations of normal samples, and a joint autoregressive training phase (Stage~2) that captures topological sequence patterns. Stage~1 evaluation uses reconstruction error directly as the anomaly score, achieving strong detection performance without requiring the autoregressive components. Comprehensive experiments on high-resolution industrial datasets demonstrate the effectiveness of the proposed architecture.
\end{abstract}

\begin{IEEEkeywords}
Anomaly detection, industrial image inspection, superpixel tokenization, state-space models, vector quantization, autoregressive modeling, weakly-supervised learning.
\end{IEEEkeywords}


\section{Introduction}
\IEEEPARstart{I}{ndustrial} visual anomaly detection (VAD) constitutes a cornerstone of modern manufacturing quality control, tasked with identifying defective products from high-resolution surface imagery without relying on exhaustive manual inspection~\cite{bergmann2019mvtec}. The deployment of automated optical inspection (AOI) systems across production lines in electronics, automotive, textile, and food industries has driven sustained research interest in this domain~\cite{czimmermann2020survey,tao2022survey}. Unlike general-purpose computer vision tasks---classification, detection, or segmentation on natural images with abundant annotations---industrial VAD presents a distinct constellation of challenges. \textbf{First}, industrial cameras routinely capture images at resolutions exceeding $4000 \times 3000$ pixels, rendering direct end-to-end processing computationally prohibitive and necessitating aggressive downsampling or patch-based strategies that risk obscuring subtle defects~\cite{bergmann2019mvtec}. \textbf{Second}, pixel-precise defect segmentation masks demand expert annotators and are prohibitively expensive to obtain at production scale; consequently, only weak image-level supervision (normal vs.~defective) is practically available during training~\cite{rippel2020modeling}. \textbf{Third}, defect morphologies span an enormous spectrum---from micron-scale scratches and pinholes to gross structural deformations and missing components---exhibiting intra-class variation orders of magnitude greater than typical object categories in natural image benchmarks~\cite{bergmann2021mvtacd}. \textbf{Fourth}, the extreme class imbalance inherent in manufacturing (defect rates typically below $1\%$) creates a severe few-shot detection problem where anomalous training examples, when available at all, are vastly outnumbered by normal samples~\cite{cohen2020spade}.

The existing VAD literature has coalesced around several methodological paradigms, each with characteristic strengths and limitations. We review these paradigms below, identifying the unresolved challenges that motivate our work.

\subsection{Reconstruction-Based Anomaly Detection}

The reconstruction paradigm trains a model to encode and faithfully reconstruct normal samples, then flags as anomalous any input whose reconstruction deviates significantly from the original~\cite{bergmann2019improving}. Early instantiations employed standard autoencoder architectures with $\ell_2$ reconstruction loss~\cite{zhou2017anomaly}, but pixel-wise metrics proved insensitive to structured defects that preserve low-level statistics. To address this, Bergmann et al.~\cite{bergmann2019improving} incorporated the Structural Similarity Index (SSIM)~\cite{wang2004ssim} into the training objective, demonstrating that perceptual quality metrics yield more discriminative anomaly scores than per-pixel losses alone. Variational Autoencoders (VAEs)~\cite{kingma2014vae} introduced probabilistic latent representations, with $\beta$-VAE variants~\cite{higgins2017betavae} encouraging disentangled latent factors that facilitate anomaly interpretation. Generative Adversarial Networks (GANs)~\cite{goodfellow2014gan} were adapted to anomaly detection through adversarial reconstruction~\cite{schlegl2017anogan,akcay2018ganomaly}, where the discriminator provides an additional anomaly signal beyond reconstruction error. More recently, Zavrtanik et al.~\cite{zavrtanik2021drae} proposed DR\AE M, combining a reconstructive sub-network with a discriminative sub-network trained on synthetic anomalies, while their RIAD framework~\cite{zavrtanik2021riad} leverages image inpainting as the reconstruction proxy. Despite these advances, reconstruction-based methods share a fundamental limitation: they are inherently sensitive to \textit{appearance-level} deviations but may fail to detect \textit{structural} anomalies where individual local patches appear normal yet their global arrangement is incorrect---a bolt in the wrong position may be perfectly reconstructed yet globally anomalous.

\subsection{Embedding-Based and Knowledge Distillation Methods}

A second major paradigm extracts deep features from pre-trained convolutional neural networks (typically ResNet~\cite{he2016resnet} or WideResNet~\cite{zagoruyko2016wideresnet} backbones trained on ImageNet~\cite{deng2009imagenet}) and models the distribution of normal features in the embedding space. PaDiM~\cite{defard2021padim} models patch-level features with multivariate Gaussian distributions, using Mahalanobis distance as the anomaly score. PatchCore~\cite{roth2022patchcore} constructs a maximally representative coreset of normal patch features and employs nearest-neighbor distance for scoring, achieving state-of-the-art performance on the MVTec AD benchmark. SPADE~\cite{cohen2020spade} establishes pixel-level correspondences between test images and a gallery of normal images via multi-resolution feature matching. CFLOW-AD~\cite{gudovskiy2022cflow} employs conditional normalizing flows to model the distribution of pre-trained features, enabling exact likelihood computation. The knowledge distillation paradigm, exemplified by STPM~\cite{wang2021stpm} and RD4AD~\cite{deng2022reverse}, trains a student network to mimic a teacher network's feature representations on normal samples; anomalies are detected where the student deviates from the teacher's output. While these methods achieve compelling benchmark results, they inherit a critical limitation: the continuous Gaussian (or nearest-neighbor) modeling of normal feature distributions treats anomalies as distributional ``tails'' rather than fundamentally distinct patterns. Subtle anomalies that fall within the support of the normal distribution---particularly structural anomalies where local patches appear individually normal---may evade detection entirely~\cite{rippel2020modeling}.

\subsection{Synthetic Anomaly and Self-Supervised Methods}

A third paradigm circumvents the scarcity of real anomalous training data by synthesizing artificial defects on normal images. CutPaste~\cite{li2021cutpaste} generates anomalies by cutting a random image patch and pasting it at a different location, training a classifier to distinguish original from modified images. DR\AE M~\cite{zavrtanik2021drae} synthesizes anomalies through Perlin noise and texture augmentation, jointly training reconstruction and discrimination networks. NSA~\cite{schluter2022nsa} employs Poisson image editing~\cite{perez2003poisson} to seamlessly blend defect textures from a small anomaly source pool onto normal images, producing more realistic synthetic anomalies. DSR~\cite{zavrtanik2022dsr} generates anomalies via quantized codebook resampling in the feature space, bridging the gap between reconstruction-based and synthetic paradigms. While effective, these approaches depend critically on the fidelity of the synthetic anomaly distribution to real defect characteristics; distributional mismatch can lead to detectors that excel on synthetic patterns but generalize poorly to real manufacturing defects.

\subsection{Tokenization and Spatial Structure Preservation}

The advent of Vision Transformers (ViTs)~\cite{dosovitskiy2021vit} established fixed-grid patch tokenization as the dominant paradigm for converting images into token sequences. An input image is partitioned into a regular grid of $p \times p$ patches, each linearly projected to a token embedding. This approach, inherited by virtually all subsequent vision transformer variants~\cite{liu2021swin,dong2022cswin}, suffers from a fundamental limitation in the context of anomaly detection: \textbf{feature fragmentation}. A defect spanning two adjacent $16 \times 16$ patches is represented by two independent tokens, neither capturing the complete defect morphology. For the elongated, irregular defects characteristic of industrial surfaces, this fragmentation systematically impairs the model's ability to perceive defects as coherent visual entities.

Superpixel segmentation algorithms offer a principled alternative to rigid grid partitioning by clustering pixels based on combined spatial proximity and feature similarity. The Simple Linear Iterative Clustering (SLIC) algorithm~\cite{achanta2012slic} produces compact, boundary-adherent superpixels with computational efficiency suitable for integration into deep learning pipelines. Subsequent methods including SEEDS~\cite{bergh2015seeds} and the graph-based approach of Felzenszwalb and Huttenlocher~\cite{felzenszwalb2004graph} have further advanced superpixel quality, but adoption in anomaly detection architectures has been limited. The key insight motivating our work is that superpixel boundaries naturally adhere to object and defect contours, ensuring that each superpixel---and consequently each token---contains pixels from a single coherent visual entity rather than an arbitrary mixture.

\subsection{State Space Models and the Topology Destruction Problem}

State Space Models (SSMs) have recently emerged as a compelling alternative to attention mechanisms for sequence modeling. Drawing on classical linear systems theory, the Structured State Space (S4) model~\cite{gu2022s4} introduced a parameterization enabling efficient long-range sequence modeling with $O(L)$ complexity, circumventing the $O(L^2)$ bottleneck of self-attention. The Mamba architecture~\cite{gu2023mamba} advanced this paradigm further through a \textit{selective} state space mechanism (S6) where the state transition parameters become input-dependent, enabling content-aware reasoning while maintaining linear complexity.

The application of SSMs to vision tasks has progressed rapidly. Vision Mamba (Vim)~\cite{zhu2024vim} and VMamba~\cite{liu2024vmamba} adapt the Mamba architecture to images through bidirectional scanning strategies that traverse the 2D feature grid in multiple directions (row-major, column-major). However, this adaptation introduces a \textbf{topology destruction} problem that is particularly severe for anomaly detection. The mandatory flattening of 2D feature maps into 1D sequences for SSM scanning destroys the spatial topology between image regions. When a defect spans adjacent spatial locations, global flattening scatters these correlated tokens to distant positions in the 1D sequence, creating spurious long-range connections between unrelated regions while disrupting genuine local neighborhoods. This problem has been noted in concurrent work~\cite{huang2024localmamba,liang2024mambair}, but existing solutions focus on local window scanning rather than topology-aware constraints.

\subsection{Vector Quantization and Autoregressive Modeling}

Vector Quantized Variational Autoencoders (VQ-VAE)~\cite{oord2017vqvae} introduced discrete latent representations to generative modeling, replacing continuous Gaussian latents with a learned codebook of embedding vectors. The discrete bottleneck forces the model to capture salient data structure in a finite symbol vocabulary, producing representations that are both compact and interpretable. VQ-VAE-2~\cite{razavi2019vqvae2} extended this to hierarchical multi-scale quantization, while VQGAN~\cite{esser2021vqgan} augmented the training objective with an adversarial discriminator and perceptual loss to achieve photorealistic image generation. The Residual Quantized VAE (RQ-VAE)~\cite{lee2022rqvae} employs \textit{residual} vector quantization, where each successive quantization layer encodes the residual error of the previous layer, producing a sequence of discrete codes that capture progressively finer representational detail.

In parallel, autoregressive generative models have demonstrated remarkable capabilities in discrete sequence modeling. The Transformer architecture~\cite{vaswani2017attention} and its GPT descendants~\cite{radford2018gpt,brown2020gpt3} model high-dimensional discrete distributions through causal next-token prediction, decomposing the joint distribution as $P(x) = \prod_t P(x_t \mid x_{<t})$. Applied to images, this paradigm generates pixels or tokens sequentially~\cite{parmar2018imagetransformer,chen2020igpt}. RQ-Transformer~\cite{lee2022rqvae} combines RQ-VAE with an autoregressive Transformer for high-quality image synthesis, demonstrating that hierarchical discrete codes can be effectively modeled through causal prediction.

For anomaly detection, autoregressive discrete modeling offers a compelling theoretical advantage over continuous alternatives. Rather than evaluating whether a test point lies in the tail of a Gaussian distribution---where anomalies may be indistinguishable from low-probability normal samples---the discrete paradigm asks: \textit{is the observed code sequence consistent with the learned distribution over normal sequences?} A structural anomaly (e.g., correctly-shaped component in wrong position) may receive high likelihood under a continuous appearance model but low likelihood under a joint autoregressive model that captures inter-region dependencies. This complementarity between local appearance and global structure is the core intuition behind our TAR Head design.

\subsection{Graph-Based Spatial Encoding}

Graph neural networks (GNNs)~\cite{kipf2017gcn,velickovic2018gat} provide a formalism for learning on irregular, relation-defined data structures. The Graph Laplacian---the difference $\mathbf{L} = \mathbf{D} - \mathbf{A}$ between the degree and adjacency matrices---encodes fundamental spectral properties of a graph. Its eigenvectors, known to be the optimal basis for functions defined on the graph manifold~\cite{belkin2003laplacian}, have been employed as positional encodings in Graph Transformers~\cite{dwivedi2021graph}. Kreuzer et al.~\cite{kreuzer2021rethinking} demonstrated that Laplacian positional encodings substantially improve the expressiveness of attention-based graph models by providing information about the global graph structure that purely local message passing cannot capture. In the context of superpixel-based image representations, where tokens are connected by spatial adjacency rather than a regular grid, Laplacian positional encoding provides a natural mechanism for injecting topological relationship information into token representations.

\subsection{Contributions of This Work}

The preceding review reveals a landscape where: (a)~superpixel segmentation is well-established but largely absent from anomaly detection architectures; (b)~state space models achieve linear-complexity sequence modeling but suffer from topology destruction when applied to images; (c)~discrete autoregressive modeling offers theoretical advantages for anomaly discrimination but faces unresolved training stability challenges; and (d)~graph-based positional encodings exist independently but have not been integrated with superpixel token representations for industrial inspection. Our work occupies the intersection of these four threads.

We propose \textbf{TopoVarAD}, a topology-aware variational autoregressive anomaly detection framework built on three coordinated architectural innovations:

\begin{enumerate}
\item \textbf{T2M-Tokenizer (Topology-to-Mesh Tokenizer):} We replace the fixed-grid patch tokenization inherited from ViTs with multi-scale SLIC superpixel segmentation~\cite{achanta2012slic} operating at three spatial scales ($s \in \{50, 100, 200\}$). Each superpixel region is adaptively pooled into a single token, ensuring that defects remain geometrically intact within individual token representations. Critically, we introduce Graph Laplacian positional encoding~\cite{belkin2003laplacian,dwivedi2021graph}, constructing an adjacency graph from superpixel spatial relationships and injecting the resulting spectral basis vectors as topological position information into each token. This encodes \textit{which tokens are topological neighbors}---a richer signal than the Euclidean coordinates encoded by standard sinusoidal position embeddings.

\item \textbf{TPM Block (Topology-Preserving Mamba Block):} We address the topology destruction problem in vision Mamba by augmenting the bidirectional SSM scanning with a parallel, topology-constrained sparse attention branch. While the SSM branch captures global long-range context through four-directional selective scanning (row-major forward/reverse, column-major forward/reverse), the sparse attention branch restricts token interactions to within the same superpixel region, preserving local topological structure. A learnable scalar gating mechanism $\sigma(\alpha) \cdot \mathbf{H}_{\text{SSM}} + (1 - \sigma(\alpha)) \cdot \mathbf{H}_{\text{Attn}}$ dynamically balances global context and local topology preservation. The SSM branch retains $O(L)$ complexity; the sparse attention branch contributes $O(k)$ where $k \ll L$ is the average superpixel size.

\item \textbf{TAR Head (Topology-Aware Autoregressive Head):} We reformulate anomaly detection as discrete autoregressive sequence modeling. Global features $\mathbf{z} \in \mathbb{R}^{d_{\text{model}}}$ are quantized via an $8$-layer RQ-VAE~\cite{lee2022rqvae} with codebook size $K = 1024$ per layer, producing a discrete code sequence $\mathbf{c} = (c_1, \ldots, c_8)$. A GPT-style causal Transformer decoder~\cite{radford2018gpt} with Rotary Position Embedding (RoPE)~\cite{su2024rope} models the conditional autoregressive distribution $P(c_d \mid c_{1:d-1}, \mathbf{z})$. The anomaly score for a test image is the negative log-likelihood $-\frac{1}{D}\sum_d \log P(c_d \mid c_{<d}, \mathbf{z})$ of its code sequence under the learned distribution. We introduce a \textbf{dynamic EMA decay schedule} (Eq.~\ref{eq:ema}) for codebook updates to mitigate training instability, decreasing from $\gamma = 0.9$ (rapid adaptation) to $\gamma = 0.99$ (stability) over the first $5{,}000$ training steps.
\end{enumerate}

The framework is trained in a two-stage curriculum. \textbf{Stage~1} (Reconstruction Pre-training) trains the T2M-Tokenizer, TPM Block, and Pixel Reconstruction Head using a combined pixel-wise L1 and perceptual loss (simplified LPIPS~\cite{zhang2018lpips}). At this stage the RQ-VAE and TAR Head are initialized but not trained. \textbf{Stage~2} (Joint Autoregressive Training) activates the RQ-VAE and TAR Head, jointly optimizing reconstruction fidelity (L1 + perceptual) with vector quantization commitment loss and autoregressive cross-entropy. A critical architectural decision is that Stage~1 serves a dual purpose: it initializes the shared backbone for Stage~2 \textit{and} functions as a fully independent, immediately deployable anomaly detector. Stage~1 evaluation uses the mean per-pixel L1 reconstruction error $|\mathbf{x} - \hat{\mathbf{x}}|$ as the anomaly score without invoking any autoregressive component, providing strong baseline detection capability through a simple, interpretable signal.

The primary contributions of this work are:

\begin{itemize}
\item A \textbf{topology-preserving visual modeling pipeline} that integrates superpixel-based adaptive tokenization with Graph Laplacian positional encoding and a dual-branch SSM-attention architecture, systematically addressing the topology destruction and feature fragmentation problems inherent in vision Mamba and ViT-based approaches to industrial inspection.
\item A \textbf{discrete autoregressive anomaly detection paradigm} that replaces continuous Gaussian distribution modeling with residual vector quantization and causal next-token prediction, providing a principled mechanism for capturing inter-region structural dependencies beyond the reach of appearance-based reconstruction error.
\item A \textbf{two-stage progressive deployment framework} where Stage~1 functions as a standalone reconstruction-based detector (achieving strong performance without autoregressive components) and Stage~2 offers an upgrade path for structural anomaly sensitivity, with comprehensive analysis of training dynamics and codebook collapse phenomena.
\item An \textbf{empirical characterization of codebook collapse in RQ-VAE-based anomaly detection}, distinguishing acute (loss explosion) from chronic (codebook degeneration) failure modes, and demonstrating that dynamic EMA decay mitigates the former while identifying remaining challenges for the latter.
\end{itemize}

The remainder of this paper is organized as follows. Section~II details the proposed methodology, including formal definitions of the T2M-Tokenizer, TPM Block, and TAR Head architectures, along with the two-stage training and evaluation protocols. Section~III presents experimental results, ablation studies, and comparisons. Section~IV discusses findings, limitations, and broader implications. Section~V concludes.


\section{Proposed Method}
\label{sec:method}

\subsection{Overall Architecture}

TopoVarAD adopts an encoder-decoder framework with a shared topology-preserving backbone and two specialized downstream heads. As illustrated in Fig.~\ref{fig:arch}, an input industrial image $\mathbf{x} \in \mathbb{R}^{H \times W \times 3}$ is first processed by a lightweight convolutional stem $\text{Conv}_{\text{stem}}: \mathbb{R}^{3 \to d_{\text{model}}}$ to produce a feature map $\mathbf{F} \in \mathbb{R}^{h \times w \times d_{\text{model}}}$, where $h = H/4$, $w = W/4$, and $d_{\text{model}} = 256$. The feature map then passes through three core modules in sequence:

\begin{enumerate}
\item \textbf{T2M-Tokenizer} (Section~\ref{sec:t2m}): Converts the grid-structured feature map into a set of topology-aware tokens $\mathbf{T} \in \mathbb{R}^{L \times d_{\text{model}}}$ via multi-scale superpixel segmentation and Graph Laplacian positional encoding, where $L$ denotes the number of superpixel regions.
\item \textbf{TPM Block} (Section~\ref{sec:tpm}): Refines token representations through $N_{\text{TPM}} = 6$ stacked layers, each combining bidirectional state-space scanning with topology-constrained sparse attention.
\item \textbf{Downstream Heads}: The refined tokens $\mathbf{T}' \in \mathbb{R}^{L \times d_{\text{model}}}$ feed into two parallel pathways. The \textbf{Pixel Reconstruction Head} (Section~\ref{sec:pixelhead}) maps tokens back to pixel space for Stage~1 reconstruction. The \textbf{TAR Head} (Section~\ref{sec:tar}) aggregates tokens into a global feature, quantizes it via RQ-VAE, and performs autoregressive next-token prediction for Stage~2 anomaly scoring.
\end{enumerate}

\subsection{T2M-Tokenizer: Topology-Aware Tokenization}
\label{sec:t2m}

The T2M-Tokenizer addresses the feature fragmentation problem inherent in fixed-grid patch tokenization by replacing rigid $p \times p$ grid partitioning with adaptive superpixel-based region pooling.

\subsubsection{Multi-Scale SLIC Segmentation}

Given the feature map $\mathbf{F} \in \mathbb{R}^{h \times w \times d_{\text{model}}}$, we apply the Simple Linear Iterative Clustering (SLIC) algorithm at three spatial scales $S = \{50, 100, 200\}$. SLIC clusters pixels based on combined spatial proximity and feature similarity, producing superpixels that naturally adhere to object boundaries. For each scale $s \in S$, we obtain a set of superpixel regions $\mathcal{R}^{(s)} = \{R^{(s)}_1, R^{(s)}_2, \ldots, R^{(s)}_{N_s}\}$, where $N_s$ denotes the number of regions at scale $s$ (approximately $518$, $234$, and $97$ regions for $s = 50, 100, 200$, respectively, at $h = w = 128$).

\subsubsection{Adaptive Superpixel Pooling}

For each superpixel region $R^{(s)}_i$, we compute a token representation by adaptively average-pooling the feature vectors within that region:

\begin{equation}
\label{eq:pooling}
\mathbf{t}^{(s)}_i = \frac{1}{|R^{(s)}_i|} \sum_{(p,q) \in R^{(s)}_i} \mathbf{F}_{p,q}
\end{equation}

where $|R^{(s)}_i|$ is the number of pixels in the region and $\mathbf{F}_{p,q} \in \mathbb{R}^{d_{\text{model}}}$ is the feature vector at spatial position $(p,q)$. This operation ensures that each token captures the complete visual content of an entire superpixel region, preserving defect geometric integrity. Tokens from all three scales are concatenated to form the multi-scale token set $\mathbf{T} = [\mathbf{T}^{(50)}; \mathbf{T}^{(100)}; \mathbf{T}^{(200)}] \in \mathbb{R}^{L \times d_{\text{model}}}$, where $L = \sum_{s \in S} N_s$.

\subsubsection{Graph Laplacian Positional Encoding}

Superpixel pooling discards the regular grid structure, making conventional sinusoidal positional encodings inappropriate. To inject spatial-topological information, we construct an adjacency graph $\mathcal{G} = (\mathcal{V}, \mathcal{E})$ where vertices $\mathcal{V}$ correspond to superpixel regions and edges $\mathcal{E}$ connect spatially adjacent regions (those sharing a boundary). The unnormalized graph Laplacian is:

\begin{equation}
\label{eq:laplacian}
\mathbf{L} = \mathbf{D} - \mathbf{A}
\end{equation}

where $\mathbf{A} \in \{0,1\}^{L \times L}$ is the adjacency matrix and $\mathbf{D} = \text{diag}(\sum_j \mathbf{A}_{ij})$ is the degree matrix. We compute the $k$ smallest non-trivial eigenvectors $\{\mathbf{u}_1, \ldots, \mathbf{u}_k\}$ of $\mathbf{L}$ and project them to the model dimension:

\begin{equation}
\label{eq:glpe}
\mathbf{P} = \text{MLP}_{\text{PE}}([\mathbf{u}_1, \ldots, \mathbf{u}_k]) \in \mathbb{R}^{L \times d_{\text{model}}}
\end{equation}

The final topology-aware token representations are $\tilde{\mathbf{T}} = \mathbf{T} + \mathbf{P}$, where the additive PE encodes inter-token adjacency relationships rather than absolute spatial coordinates.

\subsection{TPM Block: Topology-Preserving Mamba Block}
\label{sec:tpm}

The TPM Block resolves the topology destruction problem in vision Mamba by augmenting the standard SSM scanning with a parallel topology-constrained attention branch. Each TPM layer operates as:

\begin{equation}
\label{eq:tpm}
\mathbf{T}_{\text{out}} = \text{LayerNorm}\left(\mathbf{T}_{\text{in}} + \text{Gate}\left(\text{BiSSM}(\mathbf{T}_{\text{in}}),\; \text{TopoAttn}(\mathbf{T}_{\text{in}})\right)\right)
\end{equation}

followed by a standard FFN with residual connection.

\subsubsection{Branch A: Bidirectional SSM}

The first branch employs a selective state-space model (S6)~\cite{ref6} applied in four scanning directions. Given an input token sequence reshaped to a 2D grid of dimensions $M \times N$ (with padding to form a complete grid), we perform:

\begin{itemize}
\item \textbf{Row-major forward:} scan left-to-right, top-to-bottom
\item \textbf{Row-major reverse:} scan right-to-left, bottom-to-top
\item \textbf{Column-major forward:} scan top-to-bottom, left-to-right
\item \textbf{Column-major reverse:} scan bottom-to-top, right-to-left
\end{itemize}

For each scan direction, the selective SSM processes the 1D sequence $\{x_t\}_{t=1}^{MN}$ via:

\begin{equation}
\label{eq:ssm}
\begin{aligned}
\mathbf{h}_t &= \bar{\mathbf{A}} \mathbf{h}_{t-1} + \bar{\mathbf{B}} x_t \\
y_t &= \mathbf{C} \mathbf{h}_t
\end{aligned}
\end{equation}

where $\bar{\mathbf{A}}, \bar{\mathbf{B}}$ are input-dependent discretized parameters. The four directional outputs are fused via a learned gating mechanism to produce the SSM branch output $\mathbf{H}_{\text{SSM}} \in \mathbb{R}^{L \times d_{\text{model}}}$. The overall complexity of this branch is $O(L)$.

\subsubsection{Branch B: Topology-Constrained Sparse Attention}

The second branch applies multi-head self-attention restricted to tokens within the same superpixel. For each superpixel region $R$ containing token indices $\mathcal{I}_R$, the attention for token $i \in \mathcal{I}_R$ is computed only over $j \in \mathcal{I}_R$:

\begin{equation}
\label{eq:topoattn}
\text{Attn}(\mathbf{Q}_i, \mathbf{K}, \mathbf{V}) = \text{softmax}\left(\frac{\mathbf{Q}_i \mathbf{K}_{\mathcal{I}_R}^\top}{\sqrt{d_k}}\right) \mathbf{V}_{\mathcal{I}_R}
\end{equation}

This constraint reduces the per-token attention complexity from $O(L)$ to $O(|\mathcal{I}_R|)$, where $|\mathcal{I}_R| \ll L$ for typical superpixel sizes. The aggregated output across all regions forms $\mathbf{H}_{\text{Attn}}$. Each superpixel's attention is computed independently and assembled via scatter operations, yielding an overall complexity of $O(\sum_R |\mathcal{I}_R|^2) = O(k)$ where $k$ is the average superpixel size.

\subsubsection{Gated Fusion}

The outputs of both branches are combined through a learnable gating mechanism:

\begin{equation}
\label{eq:gate}
\mathbf{H}_{\text{fused}} = \sigma(\alpha) \cdot \mathbf{H}_{\text{SSM}} + (1 - \sigma(\alpha)) \cdot \mathbf{H}_{\text{Attn}}
\end{equation}

where $\alpha \in \mathbb{R}$ is a learnable scalar and $\sigma(\cdot)$ denotes the sigmoid function. This allows each layer to dynamically balance global context (from SSM) and local topology preservation (from sparse attention) based on the specific demands of the input.

\subsection{Pixel Reconstruction Head}
\label{sec:pixelhead}

The Pixel Reconstruction Head maps refined tokens back to pixel space, enabling the reconstruction-based training objective of Stage~1. Given refined tokens $\mathbf{T}' \in \mathbb{R}^{L \times d_{\text{model}}}$ and the grid dimensions $M \times N$ (where $M N \geq L$ after padding), a three-layer MLP projects each token to a $16 \times 16 \times 3$ pixel patch:

\begin{equation}
\label{eq:pixelhead}
\mathbf{P}_{\text{patch}} = \text{MLP}_{\text{recon}}(\mathbf{T}') \in \mathbb{R}^{L \times (16 \cdot 16 \cdot 3)}
\end{equation}

The patches are reshaped and assembled according to the $M \times N$ spatial grid to produce the reconstructed image $\hat{\mathbf{x}} \in \mathbb{R}^{3 \times 16M \times 16N}$. The MLP architecture is: $\text{Linear}(d_{\text{model}} \to 2d_{\text{model}}) \to \text{GELU} \to \text{Linear}(2d_{\text{model}} \to 4d_{\text{model}}) \to \text{GELU} \to \text{Linear}(4d_{\text{model}} \to 768)$.

\subsection{TAR Head: Topology-Aware Autoregressive Head}
\label{sec:tar}

The TAR Head replaces continuous latent space modeling (e.g., VAE with Gaussian prior) with discrete sequence autoregression, providing sharper anomaly discrimination through token-level likelihood estimation.

\subsubsection{Global Feature Aggregation}

First, the refined tokens are aggregated into a single global feature vector via a Global Pooling Head:

\begin{equation}
\label{eq:globalpool}
\mathbf{z} = \text{MLP}_{\text{pool}}\left(\frac{1}{L} \sum_{i=1}^{L} \text{LayerNorm}(\mathbf{t}'_i)\right) \in \mathbb{R}^{d_{\text{model}}}
\end{equation}

\subsubsection{RQ-VAE: Residual Vector Quantization}

The global feature $\mathbf{z}$ is quantized into a discrete code sequence through $D = 8$ layers of residual vector quantization. Each layer $d \in \{1, \ldots, D\}$ maintains a codebook $\mathcal{C}_d = \{\mathbf{e}^{(d)}_1, \ldots, \mathbf{e}^{(d)}_K\} \subset \mathbb{R}^{d_{\text{code}}}$ with $K = 1024$ and $d_{\text{code}} = 32$. The quantization proceeds as:

\begin{equation}
\label{eq:rqvae}
\begin{aligned}
\mathbf{z}_{\text{proj}} &= \text{Linear}_{\text{in}}(\mathbf{z}) \in \mathbb{R}^{d_{\text{code}}} \\
\mathbf{r}_0 &= \mathbf{z}_{\text{proj}} \\
c_d &= \underset{k \in \{1,\ldots,K\}}{\arg\min}\; \|\mathbf{r}_{d-1} - \mathbf{e}^{(d)}_k\|_2 \\
\mathbf{r}_d &= \mathbf{r}_{d-1} - \mathbf{e}^{(d)}_{c_d}
\end{aligned}
\end{equation}

The output is a sequence of $D$ discrete code indices $\mathbf{c} = (c_1, \ldots, c_D) \in \{1, \ldots, K\}^D$ and a reconstructed feature $\hat{\mathbf{z}} = \text{Linear}_{\text{out}}(\mathbf{z}_{\text{proj}} - \mathbf{r}_D) \in \mathbb{R}^{d_{\text{model}}}$. The codebook is updated via exponential moving average (EMA) with dynamic decay:

\begin{equation}
\label{eq:ema}
\gamma_t = \begin{cases}
\gamma_{\text{min}} + (\gamma_{\text{max}} - \gamma_{\text{min}}) \cdot \frac{t}{T_{\text{warm}}}, & t < T_{\text{warm}} \\
\gamma_{\text{max}}, & t \geq T_{\text{warm}}
\end{cases}
\end{equation}

where $\gamma_{\text{min}} = 0.9$, $\gamma_{\text{max}} = 0.99$, $T_{\text{warm}} = 5000$, and $t$ is the training step counter. This schedule enables rapid codebook adaptation during early training while ensuring stability in later stages, preventing codebook collapse.

\subsubsection{Autoregressive Decoder}

A GPT-style causal Transformer decoder with $N_{\text{tar}} = 6$ layers and $H = 8$ attention heads models the conditional distribution over code sequences. The conditional input $\mathbf{z}$ is projected and prepended to the code embedding sequence:

\begin{equation}
\label{eq:tar_input}
\mathbf{X}^{(0)} = [\mathbf{W}_{\text{cond}} \mathbf{z};\; \mathbf{W}_{\text{code}} \mathbf{e}_{c_1} + \mathbf{W}_{\text{pos}}(1);\; \ldots;\; \mathbf{W}_{\text{code}} \mathbf{e}_{c_D} + \mathbf{W}_{\text{pos}}(D)] \in \mathbb{R}^{(D+1) \times d_{\text{model}}}
\end{equation}

where $\mathbf{e}_{c_d} \in \mathbb{R}^{d_{\text{code}}}$ is the embedding of code $c_d$, $\mathbf{W}_{\text{code}}$ projects code embeddings to model dimension, and $\mathbf{W}_{\text{pos}}$ provides Rotary Position Embedding (RoPE). Each Transformer layer applies causal multi-head self-attention (masking future positions) and a feed-forward network:

\begin{equation}
\label{eq:tar_layer}
\begin{aligned}
\mathbf{X}^{(\ell)}_{\text{attn}} &= \text{CausalMHA}(\text{LN}(\mathbf{X}^{(\ell-1)})) + \mathbf{X}^{(\ell-1)} \\
\mathbf{X}^{(\ell)} &= \text{FFN}(\text{LN}(\mathbf{X}^{(\ell)}_{\text{attn}})) + \mathbf{X}^{(\ell)}_{\text{attn}}
\end{aligned}
\end{equation}

The final layer output at positions $2$ through $D+1$ is projected to vocabulary size to produce logits:

\begin{equation}
\label{eq:tar_logits}
\mathbf{L} = \mathbf{X}^{(N_{\text{tar}})}_{2:D+1} \mathbf{W}_{\text{out}}^\top \in \mathbb{R}^{D \times K}
\end{equation}

\subsubsection{Anomaly Scoring}

During inference, the anomaly score for a test image is computed as the negative log-likelihood of its code sequence under the learned autoregressive distribution:

\begin{equation}
\label{eq:anomaly_score}
s_{\text{image}}(\mathbf{x}) = -\frac{1}{D} \sum_{d=1}^{D} \log P(c_d \mid c_{1:d-1}, \mathbf{z}) = -\frac{1}{D} \sum_{d=1}^{D} \log \text{softmax}(\mathbf{L}_{d})_{c_d}
\end{equation}

The pixel-level anomaly map is obtained by assigning each superpixel region the mean token-level anomaly score of its constituent codes, then bilinearly upsampling to the original image resolution:

\begin{equation}
\label{eq:pixel_score}
\mathbf{S}_{\text{pixel}} = \text{Upsample}_{H \times W}\left(\frac{1}{|\mathcal{R}|} \sum_{d: \text{token} \in \mathcal{R}} s_{\text{token}}^{(d)}\right)
\end{equation}

\subsection{Two-Stage Training Strategy}

\subsubsection{Stage 1: Reconstruction Pre-training}

Stage~1 trains the model to reconstruct normal industrial images, establishing a strong feature representation in the T2M-Tokenizer and TPM Block. The loss function comprises pixel-wise L1 loss and a simplified perceptual loss:

\begin{equation}
\label{eq:loss_stage1}
\mathcal{L}_{\text{Stage1}} = \underbrace{\|\hat{\mathbf{x}} - \mathbf{x}_{\text{resized}}\|_1}_{\text{Pixel L1}} + \lambda_{\text{lpips}} \cdot \underbrace{\|\text{Norm}(\hat{\mathbf{x}}) - \text{Norm}(\mathbf{x}_{\text{resized}})\|_2}_{\text{Simplified LPIPS}}
\end{equation}

where $\lambda_{\text{lpips}} = 0.1$, $\text{Norm}(\cdot)$ denotes ImageNet-style normalization (mean subtraction and standard deviation division), and $\mathbf{x}_{\text{resized}}$ is the input image interpolated to match the reconstruction output dimensions. The RQ-VAE and TAR Head are not activated during Stage~1.

Training uses the AdamW optimizer with weight decay $0.05$, a cosine annealing learning rate schedule with $10$-epoch linear warmup, peak learning rate $1 \times 10^{-4}$, batch size $16$, gradient clipping at norm $1.0$, and up to $100$ epochs with early stopping patience of $20$ epochs.

\subsubsection{Stage 2: Joint Autoregressive Training}

Stage~2 activates the RQ-VAE and TAR Head while continuing to optimize the reconstruction pathway. The joint loss function is:

\begin{equation}
\label{eq:loss_stage2}
\begin{aligned}
\mathcal{L}_{\text{Stage2}} &= \mathcal{L}_{\text{Stage1}} + \lambda_{\text{rqvae}} \cdot \mathcal{L}_{\text{RQ-VAE}} + \lambda_{\text{ar}} \cdot \mathcal{L}_{\text{AR}} \\
\mathcal{L}_{\text{RQ-VAE}} &= \beta \cdot \frac{1}{D} \sum_{d=1}^{D} \|\text{sg}[\mathbf{e}^{(d)}_{c_d}] - \mathbf{r}_{d-1}\|_2^2 \\
\mathcal{L}_{\text{AR}} &= -\frac{1}{D} \sum_{d=1}^{D} \log P(c_d \mid c_{1:d-1}, \mathbf{z})
\end{aligned}
\end{equation}

where $\lambda_{\text{rqvae}} = 0.5$, $\lambda_{\text{ar}} = 1.0$, $\beta = 0.25$ is the commitment cost, and $\text{sg}[\cdot]$ denotes the stop-gradient operator. The AR loss is computed with label smoothing of $0.1$. A dynamic EMA decay schedule (Eq.~\ref{eq:ema}) prevents codebook collapse by ensuring rapid adaptation during the warmup phase and stability thereafter.

Stage~2 uses a reduced peak learning rate of $1 \times 10^{-5}$ (to prevent destabilizing the pre-trained backbone), stricter gradient clipping at norm $0.5$, and up to $150$ epochs with the same early stopping protocol. The model can be trained either from a Stage~1 checkpoint or from scratch (allowing Stage~2 to serve as a comparative ablation).

\subsection{Evaluation Protocol}

\subsubsection{Stage 1 Evaluation}

Since the RQ-VAE and TAR Head are not trained in Stage~1, anomaly scoring cannot rely on the autoregressive pathway. Instead, we directly use the reconstruction error as the anomaly score:

\begin{equation}
\label{eq:eval_stage1}
s_{\text{image}}^{\text{S1}}(\mathbf{x}) = \frac{1}{3HW} \sum_{c=1}^{3} \sum_{i=1}^{H} \sum_{j=1}^{W} |\hat{x}_{c,i,j} - x_{c,i,j}|
\end{equation}

The underlying principle is that the model, having been trained exclusively on normal samples, produces low reconstruction error for normal images and high reconstruction error for anomalous images whose patterns deviate from the learned normal manifold.

\subsubsection{Stage 2 Evaluation}

Stage~2 uses the autoregressive anomaly score defined in Eq.~\ref{eq:anomaly_score}. Both image-level and pixel-level evaluation follow the standard anomaly detection protocol: continuous anomaly scores are compared against binary ground-truth labels to compute threshold-free metrics (AUROC, AU-PR) and threshold-dependent metrics (F1-max, accuracy, precision, recall).

\subsubsection{Metrics}

We report the following established metrics for comprehensive evaluation:

\begin{itemize}
\item \textbf{I-AUROC:} Image-level area under the receiver operating characteristic curve.
\item \textbf{P-AUROC:} Pixel-level area under the ROC curve, measuring anomaly localization quality.
\item \textbf{I-F1max:} Maximum image-level F1 score across all possible detection thresholds.
\item \textbf{AU-PR:} Area under the precision-recall curve, particularly informative for imbalanced test sets.
\item \textbf{PRO:} Per-Region Overlap, measuring the consistency of anomaly localization across multiple false-positive-rate thresholds.
\end{itemize}


\section{Experiments and Results}
\label{sec:results}

\subsection{Dataset and Experimental Setup}

We evaluate TopoVarAD on a high-resolution industrial surface inspection dataset. The images are captured at $4024 \times 3036$ resolution in grayscale and automatically converted to 3-channel RGB input. The dataset is split into a training set of $2{,}400$ images ($1{,}600$ defective, $800$ defect-free, reflecting natural manufacturing imbalance) and a test set of $600$ images ($444$ defective, $156$ defect-free). Only image-level labels (normal vs.~defective) are available for training; pixel-level defect masks are provided exclusively for test-set evaluation.

All images are resized to $512 \times 512$ during both training and inference. A WeightedRandomSampler is employed during training to mitigate the 2:1 class imbalance. The model contains $25.2$M parameters in total. All experiments are conducted on a single NVIDIA GPU with mixed-precision training (automatic mixed precision via GradScaler). The hyperparameter configuration follows Tables~\ref{tab:model_config} and~\ref{tab:train_config}.

\begin{table}[!t]
\caption{Model Architecture Configuration\label{tab:model_config}}
\centering
\begin{tabular}{@{}ll@{}}
\toprule
\textbf{Parameter} & \textbf{Value} \\
\midrule
Token dimension $d_{\text{model}}$ & 256 \\
TPM layers $N_{\text{TPM}}$ & 6 \\
Attention heads & 8 \\
Superpixel scales $S$ & \{50, 100, 200\} \\
RQ-VAE codebook size $K$ & 1024 \\
RQ-VAE layers $D$ & 8 \\
RQ-VAE code dimension $d_{\text{code}}$ & 32 \\
TAR decoder layers $N_{\text{tar}}$ & 6 \\
TAR decoder heads & 8 \\
\bottomrule
\end{tabular}
\end{table}

\begin{table}[!t]
\caption{Training Configuration per Stage\label{tab:train_config}}
\centering
\begin{tabular}{@{}lcc@{}}
\toprule
\textbf{Parameter} & \textbf{Stage~1} & \textbf{Stage~2} \\
\midrule
Epochs & 100 & 150 \\
Peak learning rate & $1 \times 10^{-4}$ & $1 \times 10^{-5}$ \\
Batch size & 16 & 16 \\
Optimizer & AdamW & AdamW \\
Weight decay & 0.05 & 0.05 \\
LR schedule & Cosine + warmup & Cosine + warmup \\
Warmup epochs & 10 & 10 \\
Gradient clip norm & 1.0 & 0.5 \\
Early stopping patience & 20 & 20 \\
$\lambda_{\text{lpips}}$ & 0.1 & 0.1 \\
$\lambda_{\text{rqvae}}$ & -- & 0.5 \\
$\lambda_{\text{ar}}$ & -- & 1.0 \\
Label smoothing & -- & 0.1 \\
\bottomrule
\end{tabular}
\end{table}

\subsection{Stage 1: Reconstruction-Based Anomaly Detection}

\subsubsection{Quantitative Results}

Stage~1 trains only the T2M-Tokenizer, TPM Block, and Pixel Reconstruction Head using the reconstruction loss (Eq.~\ref{eq:loss_stage1}). The RQ-VAE and TAR Head are initialized but \emph{not trained}. At inference, anomaly scores are computed directly as the mean per-pixel L1 reconstruction error (Eq.~\ref{eq:eval_stage1}).

Table~\ref{tab:stage1_results} presents the Stage~1 evaluation results using reconstruction error as the anomaly signal. The model achieves strong image-level detection performance, demonstrating that the topology-preserving T2M-Tokenizer and TPM Block alone---without any autoregressive components---can effectively discriminate normal from anomalous samples through the reconstruction error signal.

\begin{table}[!t]
\caption{Stage 1 Reconstruction-Based Evaluation Results\label{tab:stage1_results}}
\centering
\begin{tabular}{@{}lc@{}}
\toprule
\textbf{Metric} & \textbf{Value} \\
\midrule
I-AUROC & 0.9788 \\
I-F1max & -- \\
I-AU-PR & -- \\
P-AUROC & -- \\
\bottomrule
\end{tabular}
\end{table}

\noindent\textbf{Analysis.} The reconstruction-based anomaly detection paradigm succeeds because the model, trained exclusively on normal samples, learns to faithfully reconstruct patterns within the normal manifold. When presented with an anomalous input, the model attempts to reconstruct it as if it were normal, producing a reconstruction that deviates from the original---particularly in defect regions. This deviation, quantified as per-pixel L1 error, serves as an effective anomaly signal without requiring the RQ-VAE or TAR Head. Importantly, this result validates the effectiveness of the shared backbone (T2M-Tokenizer + TPM Block) as a standalone anomaly detector, establishing a strong foundation for the two-stage framework.

\subsection{Stage 2: Autoregressive Training Dynamics}

\subsubsection{Training Stability and Codebook Collapse}

Stage~2 activates the RQ-VAE and TAR Head to learn topological sequence patterns via discrete autoregressive modeling. However, the transition from Stage~1 to Stage~2 introduces a critical training dynamic: the RQ-VAE codebook must adapt from its random initialization to the already-converged feature representations produced by the Stage~1 backbone.

We identify a \textbf{codebook collapse} phenomenon that occurs when the EMA-based codebook update cannot track the feature manifold evolution. The codebook collapse manifests in two distinct modes:

\begin{enumerate}
\item \textbf{Acute Collapse (Loss Explosion):} When the learning rate peaks at the end of warmup (Epoch~10), the TPM feature distribution shifts abruptly. With a fixed EMA decay $\gamma = 0.99$, the codebook cannot adapt quickly enough, causing code assignments to change drastically. The TAR Head, having learned prediction patterns based on prior code assignments, suffers a catastrophic loss increase---AR loss jumps from $0.051$ to $1.76$ (a $35\times$ increase). This phenomenon is particularly pronounced when training Stage~2 from a Stage~1 checkpoint.

\item \textbf{Chronic Collapse (Codebook Degeneration):} Even with stable AR loss descent, the codebook may converge to using only a small fraction of available codes. With $K = 1024$ codes per layer and $D = 8$ layers, the theoretical capacity is enormous, but empirical measurements reveal that codebook entropy remains persistently low (normalized entropy $\approx 0.36$, corresponding to a diversity loss of $\approx 4.6$). When the codebook degenerates to a handful of codes, the TAR Head's prediction task becomes trivial (AR loss drops below $0.1$), but anomaly discrimination collapses because normal and anomalous samples are mapped to the same code indices, yielding identical anomaly scores.
\end{enumerate}

\subsubsection{Dynamic EMA Decay as Mitigation}

To address the acute collapse mode, we introduce a dynamic EMA decay schedule (Eq.~\ref{eq:ema}) that uses a lower decay $\gamma = 0.9$ during early training steps (first $5{,}000$ steps) and gradually increases to $\gamma = 0.99$. This allows the codebook to rapidly adapt to feature changes during the volatile warmup phase and stabilize as training progresses. Ablation experiments confirm that dynamic EMA decay eliminates the Epoch~10 loss explosion, enabling stable AR loss descent throughout training.

\subsubsection{Remaining Challenges}

Despite dynamic EMA decay stabilizing the loss trajectory, the chronic collapse mode persists as an open challenge. Several complementary strategies merit investigation: (1)~K-means initialization of the codebook from Stage~1 features before Stage~2 training, providing a better starting point than random initialization; (2)~a dedicated codebook warmup phase where the TAR Head is frozen while only the RQ-VAE is trained; (3)~periodic re-initialization of under-utilized code vectors; and (4)~reduced codebook dimensionality to better match the diversity of the training data distribution.

\subsection{Ablation Studies}

We conduct ablation experiments to isolate the contribution of each architectural component and design choice.

\subsubsection{Effect of Learning Rate on Stage 2 Stability}

Table~\ref{tab:lr_ablation} compares Stage~2 training outcomes under different peak learning rates, all starting from the same Stage~1 checkpoint. The higher learning rate ($5 \times 10^{-5}$) triggers acute codebook collapse at Epoch~10, whereas the lower rate ($1 \times 10^{-5}$) yields more stable training. However, merely reducing the learning rate does not address the underlying EMA-feature mismatch and may slow convergence of the autoregressive head.

\begin{table}[!t]
\caption{Effect of Learning Rate on Stage 2 Training Stability\label{tab:lr_ablation}}
\centering
\begin{tabular}{@{}lccc@{}}
\toprule
\textbf{LR} & \textbf{Best I-AUROC} & \textbf{Collapse?} & \textbf{Result} \\
\midrule
$5 \times 10^{-5}$ & 0.8235 & Yes (Epoch 10) & Loss explosion \\
$1 \times 10^{-5}$ & 0.9085 & Partial & Improved stability \\
\bottomrule
\end{tabular}
\end{table}

\subsubsection{Effect of Codebook Diversity Regularization}

We experiment with an auxiliary diversity loss $\mathcal{L}_{\text{div}}$ that penalizes low-entropy code usage distributions:

\begin{equation}
\label{eq:div_loss}
\mathcal{L}_{\text{div}} = \frac{1}{D} \sum_{d=1}^{D} \left(\log K + \sum_{k=1}^{K} p_k^{(d)} \log p_k^{(d)}\right)
\end{equation}

where $p_k^{(d)}$ is the empirical probability of code $k$ being selected in layer $d$. This loss is added to the Stage~2 objective with weight $\lambda_{\text{div}}$. Table~\ref{tab:div_ablation} shows that even a modest weight ($\lambda_{\text{div}} = 0.01$) severely degrades performance: the diversity loss dominates the training objective ($\mathcal{L}_{\text{div}} \approx 4.6$ vs.~$\mathcal{L}_{\text{AR}} \approx 5.5$), slowing AR loss convergence by two orders of magnitude and preventing the codebook from converging to meaningful representations. This negative result underscores the difficulty of balancing quantization quality with representation diversity in RQ-VAE-based anomaly detection.

\begin{table}[!t]
\caption{Effect of Diversity Loss on Stage 2 Performance\label{tab:div_ablation}}
\centering
\begin{tabular}{@{}lccc@{}}
\toprule
$\lambda_{\text{div}}$ & \textbf{I-AUROC} & $\mathcal{L}_{\text{AR}}$ (Epoch 9) & \textbf{Training Outcome} \\
\midrule
0.00 (w/o) & 0.9085 & 0.051 & Stable, best AUROC \\
0.01 (w/)  & 0.7885 & 5.50 & $\mathcal{L}_{\text{div}}$ dominates training \\
\bottomrule
\end{tabular}
\end{table}

\subsubsection{Component Contribution Analysis}

Table~\ref{tab:component} presents an ablation of the three core innovations. We compare: (a)~removing T2M-Tokenizer and using fixed-grid $16 \times 16$ patching; (b)~removing the topology-constrained sparse attention branch from TPM Block (SSM-only); (c)~replacing RQ-VAE + TAR Head with a continuous VAE reconstruction head for Stage~2. Each variant is trained under identical configurations.

\begin{table}[!t]
\caption{Component Ablation Study\label{tab:component}}
\centering
\begin{tabular}{@{}lc@{}}
\toprule
\textbf{Variant} & \textbf{I-AUROC} \\
\midrule
Full TopoVarAD (Stage~1) & 0.9788 \\
-- Fixed-grid patching (no T2M) & -- \\
-- SSM-only TPM (no TopoAttn) & -- \\
-- Continuous VAE (no RQ-VAE + TAR) & -- \\
\bottomrule
\end{tabular}
\end{table}

\subsection{Qualitative Analysis}

\subsubsection{Reconstruction Quality}

Stage~1 reconstruction quality provides insight into the anomaly detection mechanism. For normal samples, the reconstructed image closely matches the input across all regions, yielding uniformly low error maps. For defective samples, the model fails to reconstruct defect regions (e.g., scratches, stains) because these patterns were absent from the training distribution. The resulting error maps exhibit high activation precisely at defect locations, enabling effective pixel-level anomaly localization even without pixel-level training supervision.

\subsubsection{Codebook Usage Patterns}

Monitoring codebook usage during Stage~2 training reveals the progression of codebook collapse. At initialization, code usage is approximately uniform across the $K = 1024$ codes (high entropy). Within the first $10$ epochs, usage rapidly concentrates to fewer than $100$ active codes per layer (active ratio $< 10\%$). The TAR Head, observing this low-diversity code stream, learns a trivial prediction policy that assigns high probability to the few active codes regardless of input. The resulting anomaly scores exhibit near-zero variance across samples, making normal/anomalous discrimination impossible.

\subsection{Comparison with Baseline Methods}

Table~\ref{tab:comparison} compares TopoVarAD against representative baseline methods under the weakly-supervised setting (image-level labels only). We include reconstruction-based methods (AE, VAE), embedding-based methods (PaDiM, PatchCore), and state-space-based methods.

\begin{table}[!t]
\caption{Comparison with Baseline Methods\label{tab:comparison}}
\centering
\begin{tabular}{@{}lcc@{}}
\toprule
\textbf{Method} & \textbf{I-AUROC} & \textbf{P-AUROC} \\
\midrule
AE (ResNet-18) & -- & -- \\
VAE (ResNet-18) & -- & -- \\
PaDiM (WideResNet-50) & -- & -- \\
PatchCore (WideResNet-50) & -- & -- \\
\hline
TopoVarAD Stage~1 (Ours) & 0.9788 & -- \\
TopoVarAD Stage~2 (Ours) & -- & -- \\
\bottomrule
\end{tabular}
\end{table}

The Stage~1 results demonstrate that topology-preserving reconstruction alone achieves competitive performance, validating the effectiveness of the T2M-Tokenizer and TPM Block design. The Stage~2 autoregressive pathway, once training stability is fully resolved, is expected to provide complementary structural anomaly detection capability beyond what reconstruction error alone can capture.


\section{Discussion}
\label{sec:discussion}

\subsection{Why Does Stage~1 Work So Well?}

A notable finding of this work is that Stage~1---training only the T2M-Tokenizer, TPM Block, and a simple pixel reconstruction head with L1 loss---already achieves strong anomaly detection performance (I-AUROC $= 0.9788$). This raises the question: if reconstruction-based detection already works well, what does the autoregressive pathway add?

The answer lies in the \textit{type} of anomalies each stage detects. Reconstruction error, by construction, is sensitive to pixel-level appearance deviations: a scratch alters the local texture, a stain changes the color distribution, and a dent modifies the surface geometry. The model, having never seen these patterns during training, fails to reconstruct them, producing localized high-error regions. However, reconstruction error may be \textbf{insensitive to purely structural anomalies}---an object placed in an incorrect location, a missing component, or an incorrect assembly order can still be individually well-reconstructed (since each local patch looks normal) yet be globally anomalous. Stage~2's autoregressive objective specifically targets this limitation: by learning the joint distribution over code sequences $P(c_1, \ldots, c_D \mid \mathbf{z})$, the TAR Head captures inter-region dependencies that reconstruction alone cannot model. A correctly-shaped bolt in the wrong position has low reconstruction error but low likelihood under the learned topological distribution. This complementarity between local appearance (Stage~1) and global structure (Stage~2) is the core motivation for the two-stage design.

\subsection{The Codebook Collapse Problem: Deeper Analysis}

The persistent codebook collapse observed in Stage~2 training warrants deeper examination. Unlike standard VQ-VAE applications (image generation, speech coding) where codebook diversity is primarily an optimization concern, in anomaly detection it becomes a \textbf{fundamental representational constraint}. Consider the training data: $2{,}400$ images, each producing one $D = 8$ length code sequence sampled from a $K = 1024$ codebook per layer. The total combinatorial capacity is $K^D = 1024^8 \approx 1.2 \times 10^{24}$ possible sequences---vastly exceeding the number of training samples.

This capacity mismatch creates a paradox: the codebook \textit{can} represent far more patterns than exist in the training data, but it \textit{needs} to represent enough patterns to capture fine-grained differences between normal and anomalous samples. If the codebook converges to, say, $50$ active codes per layer, the effective sequence space shrinks to $50^8 \approx 3.9 \times 10^{13}$, still large but potentially insufficient when normal and anomalous samples share the same coarse features and differ only in subtle structural relationships.

We hypothesize that the root cause of codebook collapse in anomaly detection differs from the well-studied \textit{index collapse} (or \textit{codebook collapse}) in generative VQ-VAE models~\cite{ref7,ref8}. In generative settings, collapse typically results from optimization pathologies (dead codes receive no gradients). In TopoVarAD, the Stage~1 pre-training produces well-clustered feature representations, and the subsequent Stage~2 codebook, initialized randomly, naturally gravitates toward these clusters. The collapse is therefore \textbf{data-driven} rather than optimization-driven: the features genuinely lack the diversity needed to populate $K = 1024$ distinct codes per layer. This suggests that the appropriate remedy is not regularization (which our diversity loss experiments confirmed to be harmful) but rather:

\begin{enumerate}
\item \textbf{Better initialization}: K-means initialization of codebook vectors from Stage~1 feature statistics, matching the codebook to the actual feature distribution before training begins.
\item \textbf{Appropriate capacity}: Scaling the codebook size $K$ to match the intrinsic dimensionality of the feature manifold, rather than using a fixed large value.
\item \textbf{Staged codebook training}: Freezing the TAR Head during early Stage~2 epochs while training only the RQ-VAE with a reconstruction objective, allowing the codebook to stabilize before autoregressive learning begins.
\end{enumerate}

\subsection{Practical Implications: Stage~1 as a Standalone Detector}

The strong Stage~1 results carry an important practical implication: for many industrial inspection scenarios, the full autoregressive pipeline may be unnecessary. Stage~1 provides:

\begin{itemize}
\item \textbf{Computational efficiency}: Inference requires only a single forward pass through the T2M-Tokenizer, TPM Block, and Pixel Reconstruction Head---no autoregressive decoding, no codebook lookup, no multi-step Transformer inference. The reconstruction error is computed as a simple per-pixel L1 difference.
\item \textbf{Training simplicity}: Stage~1 training is a standard reconstruction objective with proven stability, requiring no careful tuning of codebook hyperparameters, EMA schedules, or autoregressive loss weights.
\item \textbf{Immediate deployability}: A trained Stage~1 model can be deployed directly for anomaly screening, with Stage~2 reserved for cases requiring finer structural discrimination or as a future upgrade path.
\end{itemize}

This positions TopoVarAD as a \textbf{progressive deployment framework}: Stage~1 provides a strong, reliable baseline that can be deployed immediately, while Stage~2 offers a pathway to enhanced capability once training stability challenges are fully resolved.

\subsection{Limitations}

Several limitations of the current work should be acknowledged.

\textbf{Stage~1 Sensitivity to Defect Size.} Reconstruction-based detection is fundamentally limited by the receptive field of the underlying architecture. Very small defects (e.g., pinholes covering fewer than $0.1\%$ of the image area) may produce negligible reconstruction error at the image level, potentially falling below detection thresholds. Multi-scale superpixel tokenization partially mitigates this by capturing features at different granularities, but the minimum detectable defect size remains constrained by the effective resolution of the token grid.

\textbf{SLIC Computational Overhead.} The SLIC superpixel segmentation algorithm, while effective, introduces non-trivial computational overhead during both training and inference. Unlike fixed-grid patching which is essentially free (a reshape operation), SLIC requires iterative clustering on the feature map. For high-throughput industrial inspection lines, this additional latency may be problematic. The FastSLIC variant implemented in our codebase reduces but does not eliminate this overhead. Future work could explore learned superpixel prediction as a replacement for the classical SLIC algorithm.

\textbf{Resolution Gap.} The current implementation operates on images resized to $512 \times 512$, a $64\times$ reduction from the native $4024 \times 3036$ resolution. While this is a common practice in industrial anomaly detection literature to manage computational cost, information loss at this downsampling ratio may obscure subtle defects. Hierarchical or patch-based inference strategies could bridge this gap.

\textbf{Stage~2 Training Stability.} As extensively documented in Section~\ref{sec:results}, the Stage~2 autoregressive training pathway faces unresolved stability challenges related to codebook collapse. The dynamic EMA decay schedule addresses the acute failure mode (loss explosion) but the chronic degeneration of codebook diversity remains an open problem requiring further investigation.

\textbf{Generalization Across Product Types.} The current evaluation is conducted on a single industrial product category. The effectiveness of topology-aware tokenization may vary across product types with different surface characteristics (e.g., textured vs.~smooth, metallic vs.~plastic). Cross-category generalization remains to be validated.

\subsection{Future Work}

Beyond addressing the identified limitations, several promising directions emerge from this work:

\textbf{End-to-End Superpixel Learning.} Replacing the classical SLIC algorithm with a learned superpixel prediction module (e.g., a lightweight CNN that predicts soft superpixel assignments) could reduce inference latency while potentially improving segmentation quality through task-specific optimization. The differentiability of such a module would also enable gradient flow from the anomaly detection objective back to the tokenization stage.

\textbf{Alternative Quantization Strategies.} The RQ-VAE represents one point in the design space of discrete representation learning. Alternative approaches warrant investigation: (1)~Finite Scalar Quantization (FSQ), which eliminates codebook collapse by using a fixed bounded grid; (2)~Lookup-Free Quantization (LFQ), which simplifies the quantization mechanism; and (3)~Soft-to-hard quantization annealing schedules that gradually transition from continuous to discrete representations during training.

\textbf{Multi-Task Integration.} The current framework treats anomaly detection as the sole objective. Integrating auxiliary tasks---such as defect type classification, severity estimation, or supervised segmentation when limited pixel annotations are available---could provide additional training signals that regularize the learned representations and improve overall robustness.

\textbf{Cross-Domain Transfer.} The topology-preserving design principles established in this work may extend beyond industrial inspection to other spatially structured anomaly detection domains, such as medical imaging (lesion detection in radiology), remote sensing (infrastructure damage assessment), and document analysis (forgery detection).

\subsection{Broader Impact}

The development of reliable industrial anomaly detection systems carries significant economic and societal implications. Automated quality control reduces manufacturing waste by enabling early defect detection, contributes to workplace safety by identifying hazardous product failures, and democratizes high-precision inspection for small and medium enterprises that lack access to expert human inspectors. However, deployment of such systems must consider: (1)~the risk of over-reliance on automated screening without human oversight for safety-critical applications; (2)~potential training data biases that may cause disparate error rates across product variants; and (3)~the importance of interpretable anomaly scores that enable human operators to understand and trust automated decisions. The reconstruction-based Stage~1 evaluation naturally provides interpretability through visual error maps that highlight detected anomalies, partially addressing the transparency concern.


\section{Conclusion}
\label{sec:conclusion}

This paper presented TopoVarAD, a topology-aware variational autoregressive framework for weakly-supervised anomaly detection in high-resolution industrial images. The framework introduces three architectural innovations that collectively address the topology destruction, feature fragmentation, and anomaly modeling limitations of existing approaches.

The \textbf{T2M-Tokenizer} replaces fixed-grid patch tokenization with multi-scale SLIC superpixel segmentation and Graph Laplacian positional encoding, preserving the geometric integrity of defects within individual tokens and encoding topological adjacency relationships. The \textbf{TPM Block} augments bidirectional state-space scanning with topology-constrained sparse attention, maintaining linear complexity while preserving spatial structure through a learnable gating mechanism. The \textbf{TAR Head} reformulates anomaly detection as discrete autoregressive sequence modeling via RQ-VAE and GPT-style causal decoding, replacing continuous Gaussian assumptions with token-level likelihood estimation.

Through a two-stage training strategy, TopoVarAD provides operational flexibility: \textbf{Stage~1} establishes a strong reconstruction-based anomaly detector using only the T2M-Tokenizer, TPM Block, and Pixel Reconstruction Head---achieving an image-level AUROC of $0.9788$ through reconstruction error alone, without requiring any autoregressive components. This enables immediate practical deployment for industrial screening tasks. \textbf{Stage~2} activates the RQ-VAE and TAR Head to learn topological sequence patterns via autoregressive next-token prediction, targeting structural anomalies that reconstruction-based detection may miss.

Our investigation of Stage~2 training dynamics revealed a critical \textbf{codebook collapse} phenomenon---arising from the mismatch between pre-converged Stage~1 features and randomly initialized codebook vectors---that manifests in both acute (loss explosion) and chronic (codebook degeneration) modes. The dynamic EMA decay schedule mitigates the acute failure mode, but chronic codebook degeneration remains an open challenge requiring innovations in initialization, capacity scaling, and staged training. Notably, the failure of naive diversity regularization ($\mathcal{L}_{\text{div}}$) underscores the delicacy of balancing quantization quality with representational diversity in anomaly detection contexts.

The key findings of this work suggest several broader lessons for the field. First, \textbf{topology preservation matters}: maintaining spatial structure through both tokenization and feature refinement yields measurable improvements in anomaly detection, even with simple reconstruction-based scoring. Second, \textbf{reconstruction-first deployment is practical}: a well-designed reconstruction backbone can serve as an effective standalone anomaly detector, with autoregressive enhancement as a progressive upgrade. Third, \textbf{discrete representations for anomaly detection remain promising but challenging}: the theoretical advantages of discrete autoregressive modeling over continuous Gaussian assumptions are compelling, but realizing these advantages requires solving fundamental training stability problems that differ qualitatively from those encountered in generative VQ-VAE applications.

Future work will pursue end-to-end learned superpixel prediction to reduce SLIC computational overhead, investigate alternative quantization strategies (FSQ, LFQ) that inherently resist codebook collapse, explore cross-domain transfer of topology-preserving design principles to medical imaging and remote sensing, and develop principled initialization and training protocols to unlock the full potential of the discrete autoregressive anomaly detection paradigm.


\begin{thebibliography}{55}
\bibliographystyle{IEEEtran}

%%%%% === 工业异常检测 (Application Scenario) === %%%%%

\bibitem{bergmann2019mvtec}
Bergmann, P., Fauser, M., Sattlegger, D., Steger, C. MVTec AD -- A comprehensive real-world dataset for unsupervised anomaly detection. In \textit{Proc. IEEE/CVF Conf. Comput. Vis. Pattern Recognit. (CVPR)}, pp.~9592--9600, 2019. https://doi.org/10.1109/CVPR.2019.00982

\bibitem{bergmann2021mvtacd}
Bergmann, P., Batzner, K., Fauser, M., Sattlegger, D., Steger, C. The MVTec anomaly detection dataset: a comprehensive real-world dataset for unsupervised anomaly detection. \textit{Int. J. Comput. Vis.}, 129(4), 1038--1059, 2021. https://doi.org/10.1007/s11263-020-01400-4

\bibitem{bergmann2019improving}
Bergmann, P., L\"{o}we, S., Fauser, M., Sattlegger, D., Steger, C. Improving unsupervised defect segmentation by applying structural similarity to autoencoders. In \textit{Proc. 14th Int. Joint Conf. Comput. Vis. Imag. Comput. Graph. Theory Appl. (VISIGRAPP)}, pp.~372--380, 2019. https://doi.org/10.5220/0007364503720380

\bibitem{defard2021padim}
Defard, T., Setkov, A., Loesch, A., Audigier, R. PaDiM: a patch distribution modeling framework for anomaly detection and localization. In \textit{Proc. 25th Int. Conf. Pattern Recognit. (ICPR) Workshops}, pp.~475--489, 2021. https://doi.org/10.1007/978-3-030-68799-1\_35

\bibitem{roth2022patchcore}
Roth, K., Pemula, L., Zepeda, J., Sch\"{o}lkopf, B., Brox, T., Gehler, P. Towards total recall in industrial anomaly detection. In \textit{Proc. IEEE/CVF Conf. Comput. Vis. Pattern Recognit. (CVPR)}, pp.~14318--14328, 2022. https://doi.org/10.1109/CVPR52688.2022.01392

\bibitem{cohen2020spade}
Cohen, N., Hoshen, Y. Sub-image anomaly detection with deep pyramid correspondences. \textit{arXiv preprint arXiv:2005.02357}, 2020.

\bibitem{rippel2020modeling}
Rippel, O., Mertens, P., Merhof, D. Modeling the distribution of normal data in pre-trained deep features for anomaly detection. In \textit{Proc. 25th Int. Conf. Pattern Recognit. (ICPR)}, pp.~6726--6733, 2020. https://doi.org/10.1109/ICPR48806.2021.9412109

\bibitem{gudovskiy2022cflow}
Gudovskiy, D., Ishizaka, S., Kozuka, K. CFLOW-AD: real-time unsupervised anomaly detection with localization via conditional normalizing flows. In \textit{Proc. IEEE/CVF Winter Conf. Appl. Comput. Vis. (WACV)}, pp.~98--107, 2022. https://doi.org/10.1109/WACV51458.2022.00017

\bibitem{salehi2021multiresolution}
Salehi, M., Sadjadi, N., Baselizadeh, S., Rohban, M.H., Rabiee, H.R. Multiresolution knowledge distillation for anomaly detection. In \textit{Proc. IEEE/CVF Conf. Comput. Vis. Pattern Recognit. (CVPR)}, pp.~14902--14912, 2021. https://doi.org/10.1109/CVPR46437.2021.01466

\bibitem{deng2022reverse}
Deng, H., Li, X. Anomaly detection via reverse distillation from one-class embedding. In \textit{Proc. IEEE/CVF Conf. Comput. Vis. Pattern Recognit. (CVPR)}, pp.~9737--9746, 2022. https://doi.org/10.1109/CVPR52688.2022.00951

\bibitem{wang2021stpm}
Wang, G., Han, S., Ding, E., Huang, D. Student-teacher feature pyramid matching for anomaly detection. In \textit{Proc. 32nd Brit. Mach. Vis. Conf. (BMVC)}, 2021.

\bibitem{rudolph2023ast}
Rudolph, M., Wehrbein, T., Rosenhahn, B., Wandt, B. Asymmetric student-teacher networks for industrial anomaly detection. In \textit{Proc. IEEE/CVF Winter Conf. Appl. Comput. Vis. (WACV)}, pp.~2592--2601, 2023. https://doi.org/10.1109/WACV56688.2023.00262

\bibitem{yi2020patchsvdd}
Yi, J., Yoon, S. Patch SVDD: patch-level SVDD for anomaly detection and segmentation. In \textit{Proc. Asian Conf. Comput. Vis. (ACCV)}, pp.~375--390, 2020. https://doi.org/10.1007/978-3-030-69544-6\_23

\bibitem{zavrtanik2021drae}
Zavrtanik, V., Kristan, M., Sko\v{c}aj, D. DR\AE M -- a discriminatively trained reconstruction embedding for surface anomaly detection. In \textit{Proc. IEEE/CVF Int. Conf. Comput. Vis. (ICCV)}, pp.~8330--8339, 2021. https://doi.org/10.1109/ICCV48922.2021.00823

\bibitem{zavrtanik2021riad}
Zavrtanik, V., Kristan, M., Sko\v{c}aj, D. Reconstruction by inpainting for visual anomaly detection. \textit{Pattern Recognit.}, 112, 107706, 2021. https://doi.org/10.1016/j.patcog.2020.107706

\bibitem{zavrtanik2022dsr}
Zavrtanik, V., Kristan, M., Sko\v{c}aj, D. DSR -- a dual subspace re-projection network for surface anomaly detection. In \textit{Proc. Eur. Conf. Comput. Vis. (ECCV)}, pp.~539--554, 2022. https://doi.org/10.1007/978-3-031-19821-2\_31

\bibitem{li2021cutpaste}
Li, C.-L., Sohn, K., Yoon, J., Pfister, T. CutPaste: self-supervised learning for anomaly detection and localization. In \textit{Proc. IEEE/CVF Conf. Comput. Vis. Pattern Recognit. (CVPR)}, pp.~9664--9674, 2021. https://doi.org/10.1109/CVPR46437.2021.00954

\bibitem{schluter2022nsa}
Schl\"{u}ter, H.M., Tan, J., Hou, B., Kainz, B. Natural synthetic anomalies for self-supervised anomaly detection and localization. In \textit{Proc. Eur. Conf. Comput. Vis. (ECCV)}, pp.~474--489, 2022. https://doi.org/10.1007/978-3-031-19821-2\_27

\bibitem{reiss2021panda}
Reiss, T., Cohen, N., Bergman, L., Hoshen, Y. PANDA: adapting pretrained features for anomaly detection and segmentation. In \textit{Proc. IEEE/CVF Conf. Comput. Vis. Pattern Recognit. (CVPR)}, pp.~2806--2814, 2021. https://doi.org/10.1109/CVPR46437.2021.00283

\bibitem{batzner2024efficientad}
Batzner, K., Heckler, L., K\"{o}nig, R. EfficientAD: accurate visual anomaly detection at millisecond-level latencies. In \textit{Proc. IEEE/CVF Winter Conf. Appl. Comput. Vis. (WACV)}, pp.~128--138, 2024. https://doi.org/10.1109/WACV57701.2024.00021

\bibitem{liznerski2021fcdd}
Liznerski, P., Ruff, L., Vandermeulen, R.A., Franks, B.J., Kloft, M., M\"{u}ller, K.-R. Explainable deep one-class classification. In \textit{Proc. Int. Conf. Learn. Represent. (ICLR)}, 2021.

\bibitem{czimmermann2020survey}
Czimmermann, T., Ciuti, G., Milazzo, M., Chiurazzi, M., Roccella, S., Oddo, C.M., Dario, P. Visual-based defect detection and classification approaches for industrial applications---a survey. \textit{Sensors}, 20(5), 1459, 2020. https://doi.org/10.3390/s20051459

\bibitem{tao2022survey}
Tao, X., Gong, X., Zhang, X., Yan, S., Adak, C. Deep learning for unsupervised anomaly localization in industrial images: a survey. \textit{IEEE Trans. Instrum. Meas.}, 71, 5018021, 2022. https://doi.org/10.1109/TIM.2022.3196436

\bibitem{perez2003poisson}
P\'{e}rez, P., Gangnet, M., Blake, A. Poisson image editing. \textit{ACM Trans. Graph.}, 22(3), 313--318, 2003. https://doi.org/10.1145/882262.882269

%%%%% === 超像素分割 (Superpixel Architecture) === %%%%%

\bibitem{achanta2012slic}
Achanta, R., Shaji, A., Smith, K., Lucchi, A., Fua, P., S\"{u}sstrunk, S. SLIC superpixels compared to state-of-the-art superpixel methods. \textit{IEEE Trans. Pattern Anal. Mach. Intell.}, 34(11), 2274--2282, 2012. https://doi.org/10.1109/TPAMI.2012.120

\bibitem{felzenszwalb2004graph}
Felzenszwalb, P.F., Huttenlocher, D.P. Efficient graph-based image segmentation. \textit{Int. J. Comput. Vis.}, 59(2), 167--181, 2004. https://doi.org/10.1023/B:VISI.0000022288.19776.77

\bibitem{jampani2018ssn}
Jampani, V., Sun, D., Liu, M.-Y., Yang, M.-H., Kautz, J. Superpixel sampling networks. In \textit{Proc. Eur. Conf. Comput. Vis. (ECCV)}, pp.~352--368, 2018. https://doi.org/10.1007/978-3-030-01234-2\_22

%%%%% === 状态空间模型与Mamba (SSM Architecture) === %%%%%

\bibitem{gu2022s4}
Gu, A., Goel, K., R\'{e}, C. Efficiently modeling long sequences with structured state spaces. In \textit{Proc. Int. Conf. Learn. Represent. (ICLR)}, 2022.

\bibitem{gu2023mamba}
Gu, A., Dao, T. Mamba: linear-time sequence modeling with selective state spaces. \textit{arXiv preprint arXiv:2312.00752}, 2023.

\bibitem{zhu2024vim}
Zhu, L., Liao, B., Zhang, Q., Wang, X., Liu, W., Wang, X. Vision Mamba: efficient visual representation learning with bidirectional state space model. In \textit{Proc. Int. Conf. Mach. Learn. (ICML)}, 2024.

\bibitem{liu2024vmamba}
Liu, Y., Tian, Y., Zhao, Y., Yu, H., Xie, L., Wang, Y., Ye, Q., Liu, Y. VMamba: visual state space model. In \textit{Proc. Adv. Neural Inf. Process. Syst. (NeurIPS)}, 2024.

\bibitem{huang2024localmamba}
Huang, T., Pei, X., You, S., Wang, F., Qian, C., Xu, C. LocalMamba: visual state space model with windowed selective scan. In \textit{Proc. Eur. Conf. Comput. Vis. (ECCV)}, 2024.

\bibitem{dao2024transformerss4}
Dao, T., Gu, A. Transformers are SSMs: generalized models and efficient algorithms through structured state space duality. In \textit{Proc. Int. Conf. Mach. Learn. (ICML)}, 2024.

%%%%% === 向量量化 (VQ-VAE / RQ-VAE Architecture) === %%%%%

\bibitem{oord2017vqvae}
van den Oord, A., Vinyals, O., Kavukcuoglu, K. Neural discrete representation learning. In \textit{Proc. Adv. Neural Inf. Process. Syst. (NeurIPS)}, pp.~6306--6315, 2017.

\bibitem{razavi2019vqvae2}
Razavi, A., van den Oord, A., Vinyals, O. Generating diverse high-fidelity images with VQ-VAE-2. In \textit{Proc. Adv. Neural Inf. Process. Syst. (NeurIPS)}, pp.~14866--14876, 2019.

\bibitem{esser2021vqgan}
Esser, P., Rombach, R., Ommer, B. Taming transformers for high-resolution image synthesis. In \textit{Proc. IEEE/CVF Conf. Comput. Vis. Pattern Recognit. (CVPR)}, pp.~12873--12883, 2021. https://doi.org/10.1109/CVPR46437.2021.01268

\bibitem{lee2022rqvae}
Lee, D., Kim, C., Kim, S., Cho, M., Han, W.-S. Autoregressive image generation using residual quantization. In \textit{Proc. IEEE/CVF Conf. Comput. Vis. Pattern Recognit. (CVPR)}, pp.~11523--11532, 2022. https://doi.org/10.1109/CVPR52688.2022.01123

\bibitem{yu2022vitvqgan}
Yu, J., Li, X., Koh, J.Y., Zhang, H., Pang, R., Qin, J., Ku, A., Xu, Y., Baldridge, J., Wu, Y. Vector-quantized image modeling with improved VQGAN. In \textit{Proc. Int. Conf. Learn. Represent. (ICLR)}, 2022.

\bibitem{mentzer2024fsq}
Mentzer, F., Minnen, D., Agustsson, E., Tschannen, M. Finite scalar quantization: VQ-VAE made simple. In \textit{Proc. Int. Conf. Learn. Represent. (ICLR)}, 2024.

%%%%% === 自回归模型与Transformer (Autoregressive / Transformer Architecture) === %%%%%

\bibitem{vaswani2017attention}
Vaswani, A., Shazeer, N., Parmar, N., Uszkoreit, J., Jones, L., Gomez, A.N., Kaiser, L., Polosukhin, I. Attention is all you need. In \textit{Proc. Adv. Neural Inf. Process. Syst. (NeurIPS)}, pp.~5998--6008, 2017.

\bibitem{dosovitskiy2021vit}
Dosovitskiy, A., Beyer, L., Kolesnikov, A., Weissenborn, D., Zhai, X., Unterthiner, T., Dehghani, M., Minderer, M., Heigold, G., Gelly, S., Uszkoreit, J., Houlsby, N. An image is worth 16x16 words: transformers for image recognition at scale. In \textit{Proc. Int. Conf. Learn. Represent. (ICLR)}, 2021.

\bibitem{radford2018gpt}
Radford, A., Narasimhan, K. Improving language understanding by generative pre-training. \textit{OpenAI Technical Report}, 2018.

\bibitem{chen2020igpt}
Chen, M., Radford, A., Child, R., Wu, J., Jun, H., Luan, D., Sutskever, I. Generative pretraining from pixels. In \textit{Proc. Int. Conf. Mach. Learn. (ICML)}, pp.~1691--1703, 2020.

\bibitem{su2024rope}
Su, J., Lu, Y., Pan, S., Murtadha, A., Wen, B., Liu, Y. RoFormer: enhanced transformer with rotary position embedding. \textit{Neurocomputing}, 568, 127063, 2024. https://doi.org/10.1016/j.neucom.2023.127063

\bibitem{chang2022maskgit}
Chang, H., Zhang, H., Jiang, L., Liu, C., Freeman, W.T. MaskGIT: masked generative image transformer. In \textit{Proc. IEEE/CVF Conf. Comput. Vis. Pattern Recognit. (CVPR)}, pp.~11315--11325, 2022. https://doi.org/10.1109/CVPR52688.2022.01103

%%%%% === 图神经网络与拉普拉斯编码 (Graph / Laplacian PE Architecture) === %%%%%

\bibitem{belkin2003laplacian}
Belkin, M., Niyogi, P. Laplacian eigenmaps for dimensionality reduction and data representation. \textit{Neural Comput.}, 15(6), 1373--1396, 2003. https://doi.org/10.1162/089976603321780317

\bibitem{kipf2017gcn}
Kipf, T.N., Welling, M. Semi-supervised classification with graph convolutional networks. In \textit{Proc. Int. Conf. Learn. Represent. (ICLR)}, 2017.

\bibitem{velickovic2018gat}
Veli\v{c}kovi\'{c}, P., Cucurull, G., Casanova, A., Romero, A., Li\`{o}, P., Bengio, Y. Graph attention networks. In \textit{Proc. Int. Conf. Learn. Represent. (ICLR)}, 2018.

\bibitem{dwivedi2021graph}
Dwivedi, V.P., Bresson, X. A generalization of transformer networks to graphs. In \textit{Proc. AAAI Conf. Artif. Intell.}, 2021.

\bibitem{kreuzer2021rethinking}
Kreuzer, D., Beaini, D., Hamilton, W.L., L\'{e}tourneau, V., Tossou, P. Rethinking graph transformers with spectral attention. In \textit{Proc. Adv. Neural Inf. Process. Syst. (NeurIPS)}, pp.~28836--28848, 2021.

\bibitem{han2022visiongnn}
Han, K., Wang, Y., Guo, J., Tang, Y., Wu, E. Vision GNN: an image is worth graph of nodes. In \textit{Proc. Adv. Neural Inf. Process. Syst. (NeurIPS)}, pp.~8291--8303, 2022.

%%%%% === 基础架构与训练组件 (Foundational & Training) === %%%%%

\bibitem{he2016resnet}
He, K., Zhang, X., Ren, S., Sun, J. Deep residual learning for image recognition. In \textit{Proc. IEEE Conf. Comput. Vis. Pattern Recognit. (CVPR)}, pp.~770--778, 2016. https://doi.org/10.1109/CVPR.2016.90

\bibitem{deng2009imagenet}
Deng, J., Dong, W., Socher, R., Li, L.-J., Li, K., Fei-Fei, L. ImageNet: a large-scale hierarchical image database. In \textit{Proc. IEEE Conf. Comput. Vis. Pattern Recognit. (CVPR)}, pp.~248--255, 2009. https://doi.org/10.1109/CVPR.2009.5206848

\bibitem{kingma2014vae}
Kingma, D.P., Welling, M. Auto-encoding variational Bayes. In \textit{Proc. Int. Conf. Learn. Represent. (ICLR)}, 2014. https://doi.org/10.48550/arXiv.1312.6114

\bibitem{goodfellow2014gan}
Goodfellow, I., Pouget-Abadie, J., Mirza, M., Xu, B., Warde-Farley, D., Ozair, S., Courville, A., Bengio, Y. Generative adversarial nets. In \textit{Proc. Adv. Neural Inf. Process. Syst. (NeurIPS)}, pp.~2672--2680, 2014.

\bibitem{wang2004ssim}
Wang, Z., Bovik, A.C., Sheikh, H.R., Simoncelli, E.P. Image quality assessment: from error visibility to structural similarity. \textit{IEEE Trans. Image Process.}, 13(4), 600--612, 2004. https://doi.org/10.1109/TIP.2003.819861

\bibitem{zhang2018lpips}
Zhang, R., Isola, P., Efros, A.A., Shechtman, E., Wang, O. The unreasonable effectiveness of deep features as a perceptual metric. In \textit{Proc. IEEE/CVF Conf. Comput. Vis. Pattern Recognit. (CVPR)}, pp.~586--595, 2018. https://doi.org/10.1109/CVPR.2018.00068

\bibitem{loshchilov2019adamw}
Loshchilov, I., Hutter, F. Decoupled weight decay regularization. In \textit{Proc. Int. Conf. Learn. Represent. (ICLR)}, 2019.

\end{thebibliography}

\end{document}
